\documentclass[abstracton,parskip=half,twoside,open=right]{scrreprt}
\usepackage[ngerman,english]{babel}
\usepackage[utf8]{inputenc}
\usepackage[T1]{fontenc}
\usepackage{graphicx}
\usepackage{hyperref}
\usepackage{todonotes}
\usepackage{listings}
\usepackage{lstautogobble}
\usepackage{url}
\usepackage{amssymb}
\usepackage{amsthm}
\usepackage{textcomp}

% Set the default settings of Komascript for oneside use. With 1cm for the
% binding, this will result in 2cm inner margin; and 3cm outer margin. This
% is less bad then the default 2 left, 4 right IMO (with binding correction
% 3 left, 4 right).
\usepackage[left=3cm,right=3cm,bottom=4cm,includefoot,top=3cm]{geometry}
\DeclareUnicodeCharacter{03BC}{\textmu}
% Setup the title page. Abuse some fields to get a nice layout
\author{Julian Andres Klode\\Altenritter Str. 4\\34270 Schauenburg\\Matriculation number: 2403668\\\\Supervised by: Prof. Klaus Ostermann,\\\ Yufei Cai}
\subject{\includegraphics[height=.13\textheight]{graphics/Logo_uni}\\Programming Languages and Software Technology Group\\Prof. Ostermann\\[1.5cm]}
\title{Performance and Interfaces of Datatype-Generic Scala Programs}
\subtitle{Bachelor Thesis}
\date{\today}

% Make sure that a bibliography entry does not have a page break
\usepackage{etoolbox}
\apptocmd{\thebibliography}{\interlinepenalty 10000\relax}{}{}

% Setup PDF support
\hypersetup{
   hidelinks=true,
   pdfauthor={Julian Andres Klode},
   pdftitle={Bachelor Thesis: Performance and Interfaces of Datatype-Generic Scala Programs},
   pdfsubject={Programming Languages},
}

% Setup listings
\lstdefinelanguage{scala}{
  keywords={%
          abstract,case,catch,class,def,do,else,extends,%
          false,final,finally,for,forSome,if,implicit,import,lazy,%
          match,new,null,object,override,package,private,protected,%
          return,sealed,super,this,throw,trait,true,try,type,%
          val,var,while,with,yield},
  otherkeywords={=>,<-,<\%,<:,>:,\#,@},
  sensitive=true,
  morecomment=[l]{//},
  morecomment=[n]{/*}{*/},
  morestring=[b]",
  morestring=[b]"""
}[keywords,comments,strings]

\lstset{
      basicstyle=\footnotesize,% Not sure yet. Code is large.
      frame=tb,
      language=scala,
      autogobble=true,
      basewidth={0.5em,0.45em},
      literate={=}{{$=$}}1
               {>}{{$>$}}1 {<}{{$<$}}1 {\\}{{$\lambda$}}1
               {\\\\}{{\char`\\\char`\\}}1
               {->}{{$\rightarrow$}}2 {>=}{{$\geq$}}2
               {←}{{$\leftarrow$}}2 {<-}{{$\leftarrow$}}2
               {⇒}{{$\Rightarrow$}}2 {=>}{{$\Rightarrow$}}2
               {<=}{{$\leq$}}2 {|}{{$\mid$}}1
               {μ}{{\textmu}}1
}

% A style that breaks lines
\lstdefinestyle{breaklines}{
      breaklines=true,
      breakatwhitespace=true,
      prebreak=\raisebox{0ex}[0ex][0ex]{\ensuremath{\hookleftarrow}},
      postbreak=\raisebox{0ex}[0ex][0ex]{\ensuremath{\hookrightarrow\space}}
}

% Add the current section to the listing number
\AtBeginDocument{\renewcommand{\thelstlisting}{\thesection.\arabic{lstlisting}}}
\makeatletter
\def\l@lstlisting#1#2{\@dottedtocline{1}{1.5em}{3.5em}{#1}{#2}}
\makeatother

% Set tocdepth to 1, otherwise we get a multi-page table of contents
\setcounter{tocdepth}{1}

% Additional commands
\newcommand{\cd}[1]{\lstinline[basicstyle=\small\sffamily]$#1$}
\newcommand{\summary}[1]{\textit{\textbf{In short:} #1}}

% An example environment
\theoremstyle{definition}
\newtheorem{example}{Example}[subsection]

\begin{document}

    % Title page, abstract, toc: Roman page numbers starting at i
    \pagenumbering{roman}

    \maketitle{}
    \begingroup
\let\titlepage\par
\let\endtitlepage
\let
\selectlanguage{ngerman}
\begin{abstract}
Generische Programmierung ist eine Vorgehensweise, die es ermöglicht, Funktionen
über die Struktur eines Datentypes zu definieren, anstelle von eines spezifischen
Datentypes. Dies erlaubt es einer Funktion auf beliebigen Datentypen zu operieren,
selbst auf solchen, die dem Programmier der Funktion noch nicht bekannt waren.

In dieser Bachelorarbeit werde ich Portierungen von drei Haskell Bibliotheken
(LIGD, EMGM, Uniplate) nach Scala vorstellen und ihre Funktionalität und Geschwindigkeit miteinander
sowie mit Shapeless vergleichen.

Ich werde außerdem zeigen, dass LIGD in Scala erweitert werden kann, sodass sie
beinahe so mächtig wie EMGM ist, indem ich das Problem der Erweiterbarkeit
von Funktionen löse, welches bisherige Implementierungen stark beeinträchtige.
\end{abstract}
\selectlanguage{english}
\begin{abstract}
Generic programming is a technique that allows functions to be defined over
the shape of an type, rather than a specific type; thus allowing one function
to be used on objects of arbitrary types, even those that were not known to
the programmer of the function.

In this thesis, I will present ports of three Haskell libraries (LIGD, EMGM, Uniplate) to Scala,
and compare their features and performance against each other and
against shapeless, a native Scala library.

I will show that LIGD in Scala can be extended to be almost as powerful as EMGM;
solving the issue with extensibility that made it vastly inferior in previous
implementations.


\end{abstract}
\endgroup

    \selectlanguage{english}
    \tableofcontents{}
    \cleardoublepage

    % Main part: Use arabic page numbers starting at 1
    \pagenumbering{arabic}
    \setcounter{page}{1}

    \chapter{Introduction}
Datatype-generic operations are operations that are defined over the shapes
of objects rather than any specific types of objects. This is different from
the concept of parametric polymorphism, also called generic programming in
languages like Scala:
Where parametric polymorphism abstracts away the elements of a container,
datatype-generic programming essentially abstracts away the container.

In the rest of this thesis, `generic programming' means `datatype-generic
programming'.


\begin{example}[Parametric polymorphism]
The type constructor \cd{List} is an example for parametric polymorphism. For
example, we can generate the length of a list:

\begin{lstlisting}
    def length[T](list: List[T]) = ...
\end{lstlisting}
\end{example}

\begin{example}[Datatype-generic operation]
A datatype-generic operation is an operation that abstracts the outer type, and
works with specific inner types. For example, a sum that can work on containers
of integers; that is, something like:
\begin{lstlisting}
    def sum[C[_]](list: C[Int]) = ...
\end{lstlisting}
\end{example}

There are two basic ways of how datatype-generic programming libraries expose
their interfaces:

\begin{enumerate}
    \item By representing objects as sums and products
    \item By providing combinators for performing operations
\end{enumerate}

In the first style, a generic operation will traverse the data structure and
perform its operation. In the second style, the generic library takes care of
the traversal and offers some combinators to allow working with the elements.
The first style is thus similar to traversing a list manually whereas
the second style resembles a fold.

\section*{What defines a good generic programming library?}
In `Comparing Libraries for Generic Programming in Haskell'~\cite{DBLP:conf/haskell/RodriguezJJGKO08}
, three basic types of operations a library should support were identified.
These are (where \cd{T} is a fixed type):

\begin{enumerate}
    \item \textbf{Consumers} take a generic object and return an object of a fixed type, they have
           a type of the form \cd{a => T}

           Example: The sum function from above.

    \item \textbf{Transformers} map from one generic type to another or the same, they have a type of the form
            \cd{a => a} or \cd{a => b}

            Example: A function that increments every integer in an object, returning the modified one.

    \item \textbf{Producers} produce a generic object, they have a type of the form: \cd{T => a}

            Example: A function that takes an integer and the representation of
                     a tree and returns a tree.
\end{enumerate}


In addition to supporting the three type of operations above, a good generic
programming library should also fulfil the following criteria:

\paragraph{Ad-hoc cases:} Can a function deal with a specific type in an ad-hoc way?
\\
For example, if the library exposes sums and products, can a function
only work on sums and products, or can it match some types directly? That is,
can we write a function that collects all objects of a specific type?

\paragraph{Extensibility:} Can an existing function be extended with new ad-hoc cases?
\\
For example, imagine a pre-existing function to calculate the depth of a data
structure. If we write a new tree type that stores the depth in its nodes, can
we extend the pre-existing function to simply return that value?

\paragraph{Multi-parametric functions:} Can there be functions with two generic parameters?
\\
For example, we might want to define a generic equality function, a generic
comparison function, or a function that adds generic values. One other question
to consider here is whether the parameters can be of different types.

\paragraph{First-class generics:} Can a generic function be passed to another function?
\\
For example, suppose a function \cd{gmapQ} that takes another generic function \cd{f}
and an object and maps the function \cd{f} over all members of the object, constructing
a list of new objects.

\paragraph{Automatic representations:} Generic libraries are usually implemented
using some sort of representation objects that describe the shape of
objects. Is it possible to create those automatically?

\paragraph{Separate compilation:} Can we introduce a new type in a new module
and use it with a generic function in an existing module?
\textit{(All of the libraries in this thesis support this)}


\paragraph{Constructor names:} If objects are represented as sums of products,
can those products have names?
For example, a pretty printer would want to print constructor names.


\section*{Current state}
Generic programming is primarily used by the Haskell community. There are
several libraries available and the Glasgow Haskell Compiler already provides
a library as part of its \cd{base} package, the \cd{Data.Data} module. This
module is based on another library called `Scrap your boilerplate'\cite{DBLP:conf/tldi/LammelJ03}
and is combinator-based.

Scala has one generic programming library called `Shapeless'. Shapeless was
initially derived from scrap your boilerplate, but became more sophisticated
over the time, adding features like heterogeneous lists. This is covered
in more detail in section~\ref{shapeless}.

    \chapter{Porting from Haskell to Scala}

This chapter introduces some basic concepts for translating Haskell code
to Scala. Most of these concepts were explained by Oliveira and Gibbons~\cite{scalagp}, this
chapter summarizes them.

\section{Type classes and Instances}

Haskell as a concept of type classes and instances, whereas Scala offers
traits. Traits differ from type classes in that their methods always contain
an implicit \cd{this} object, that is, they combine data and functions;
whereas type classes do not reference any data.

Given a type class like
\begin{lstlisting}[language=Haskell]
  class MyEq a where
    myEquals :: a -> a -> Boolean
    myEquals a b = False
\end{lstlisting}
how do we translate this to Scala? There are two approaches:

\subsection{Object-oriented translation}
The object-oriented approach uses traits just like interfaces in
Java are used. For example, the \cd{MyEq} type class would be translated
to the following trait:
\begin{lstlisting}
  trait MyEq[T] {
    def myEquals(b: T) : Boolean = false
  }
\end{lstlisting}
In order to implement this trait for a class the class must extend the
trait and override the methods of the trait:
\begin{lstlisting}
  class MyInt(val i: Int) extends MyEq[MyInt] {
    override def myEquals(b: MyInt) : Boolean = i == b.i
  }
\end{lstlisting}
This already shows one deficiency of the traits approach: We cannot
implement a trait for a type without modifying that type itself.

Writing a function using this approach is straight forward:
\begin{lstlisting}
  def myEquals[T <: MyEq[T]](a: T, b: T) = a.myEquals(b)
\end{lstlisting}

This style makes it easy to specify contexts, so for example, if we have
a class \cd{MyOrd} for ordering that also requires equality, we can write
\begin{lstlisting}
  trait MyOrd[T] extends MyEq[T] {
    def myLessOrEqual(b: T) : Boolean
  }
\end{lstlisting}

If we want to translate a type class containing a method that does not take a
parameter of the type A, this style of writing becomes inconsistent.

\begin{example} Consider the type class
\begin{lstlisting}[language=Haskell]
  class PerformerFactory a where
    factor :: a
    perform :: a -> Boolean
\end{lstlisting}
We can only directly translate the \cd{perform} method to Scala. The \cd{factor}
method can only be translated by use of a generic function -- this will
require a stub implementation, that instances can specialize.
\begin{lstlisting}
  trait PerformerFactory[A] {
    def factor : A
  }
  def perform[A](a: A) : Boolean = throw new Exception
\end{lstlisting}
The end result is inconsistent, because one method is now a method of
instance objects whereas the other is translated to a global function. It
also permits instances of the trait that only implement the \cd{factor}
method.
\end{example}
\subsection{The functional approach}
Remember the type class MyEq:
\begin{lstlisting}[language=Haskell]
  class MyEq a where
    myEquals :: a -> a -> Boolean
    myEquals a b = False
\end{lstlisting}
We can ignore the implicit this parameter in traits and translate the type
class as:
\begin{lstlisting}
  trait MyEq[A] {
    def myEquals(a: A, b: A) = false
  }
  def myEquals[T](a: T, b: T)(implicit eq: MyEq[T]) = eq.myEquals(a, b)
\end{lstlisting}
We also provided a global \cd{myEquals} function to make our life easier.

Now in order to implement the trait for a type, we need to define an instance
of the trait for this type. Instances of those traits correspond to dictionaries
in Haskell, and can be implemented like this:
\begin{lstlisting}
  implicit def MyEqInt = new MyEq[Int] {
    override def myEquals(a: Int, b: Int) = a == b
  }
\end{lstlisting}
The \cd{implicit} keyword is a special syntactic sugar provided by Scala
that makes our work easier. Without it, if we wanted to compare two objects,
we would need to manually pass the instance of the trait to the function,
like this:
\begin{lstlisting}
  > myEquals(1, 1)(MyEq)
  true
\end{lstlisting}
By making \cd{MyEqInt} and the \cd{eq} parameter of \cd{myEquals} implicit,
the Scala compiler can automatically figure out that \cd{MyEqInt} is the
correct argument for the \cd{eq} parameter, and we can now run:
\begin{lstlisting}
  > myEquals(1, 1)
  true
\end{lstlisting}


We can also implement instance that require some sort of context. For
example, if we want to implement \cd{MyEq} for tuples of objects implementing
\cd{MyEq} as well, we can write:
\begin{lstlisting}
  implicit def myEqTuple[A,B](implicit ea: MyEq[A], eb: MyEq[B]) =
    new MyEq[(A,B)]{
    override def myEquals(a: (A,B), b: (A,B)) = ea.myEquals(a._1, b._1)
                                    && eb.myEquals(a._2, b._2)
  }
\end{lstlisting}
Now we can compare tuples: For example, in the following example, Scala will
automatically infer the correct instance of \cd{MyEq[(Int, Int)]}:
\begin{lstlisting}
  > myEquals((1,2), (1,2))
  true
\end{lstlisting}
We can still pass  it manually though
\begin{lstlisting}
  > myEquals((1,2), (1,2))(myEqTuple(MyEqInt, MyEqInt))
  true
\end{lstlisting}

\section{Making the most of type inference}
\label{type-inference}

Compared to Haskell, Scala's type inference is weaker. For example, when
defining a function, the type of parameters need to be defined. An exception
to this are anonymous functions. If the Scala compiler knows which type the
function must have, it can automatically infer the type of an anonymous
function. Thus, in order to avoid repetitious type information in the code,
we can change our \cd{MyEq} example to:
\begin{lstlisting}
  trait MyEq[A] {
    def myEquals : (A => A => Boolean) = a => b => false
  }
  implicit def MyEqInt = new MyEq[Int] {
    override def myEquals = a => b => a == b
  }
\end{lstlisting}
Type annotations are not required in method implementations anymore. This
technique is shown in the port of the EMGM library.

% Discovered: 2014-08-04
Another thing that seems to matter a lot is the order of parameters. If
the less complex types are left of the more complex types, the more
complex types can be inferred more easily. For example, in Haskell,
\cd{foldl} has the signature:
\begin{lstlisting}[language=Haskell]
  foldl :: (a -> b -> a) -> a -> [b] -> a
\end{lstlisting}
Scala does not seem to be able to understand this, it will fail to
infer the type of the first argument. Moving the first argument to
the end fixes the issue, as Scala can then simply infer from left
to right. It basically looks like this in Scala (if you ignore that
the first argument is actually the object which foldLeft is a method
of):
\begin{lstlisting}
  foldLeft[a,b] : [b] => a => (a => b => a) => a
\end{lstlisting}

    \chapter{Libraries}
\label{libraries}

Haskell provides many libraries to assist with generic programming. Scala on
the other hand only has the Shapeless library. In this chapter, we will
translate some Haskell libraries to Scala and in some cases extend them with
new functionality made possible by the Scala type system.

\section{LIGD$^{+}$}
\summary{The straight-forward solution to most problems}
\lstinputlisting[firstline=53,lastline=68,float=ht,caption=LIGD Basics,label=ligd-base]{scala/src/ligd.scala}

LIGD\cite{Cheney:2002:LIG:581690.581698} is a simple library that provides a \cd{Rep} trait that describes
the structure of a type using sums, products, and scalar values. This
version of LIGD is more powerful than the original one, though, as we
will see later.

Listing~\ref{ligd-base} shows the basic \cd{Rep} trait with instances for
scalar types, sums (represented using \cd{Either}), and products (represented
using \cd{Tuple2}). As explained in the previous chapter, we use implicit
here to allow the compiler to automatically infer the correct \cd{Rep}
instance for generic function calls.


\subsection{User-defined types}
In order to work with user-defined (case) classes, those need to be
translated to a representation using sums and products. We need an
isomorphism between the user-defined class and a sum of product type. First
of all, we can define an isomorphismus as:
\begin{lstlisting}[gobble=2]
  class EP[B, C](val from: B ⇒ C, val to: C ⇒ B)
\end{lstlisting}
Now we can define a \cd{Rep} instance, that represents values of a custom
type \cd{B} and converts it to a sum of product notation of type \cd{C},
something like this:
\begin{lstlisting}[gobble=2]
  case class RType[C, B](val c: Rep[C], val ep: EP[B, C]) extends Rep[B]
\end{lstlisting}
There is one problem though: The representation of \cd{C} may be recursive,
for example a list type is a sum of nothing, and a product of an element and
the list type. Thus, the parameter \cd{c} must be lazy. Scala does not allow
lazy \cd{val}s on case classes, though, so this needs to be worked around.

There are two approaches: One is encoding the thunk manually, like:
\begin{lstlisting}[gobble=2]
  case class RType[C, B](val c: () => Rep[C], val ep: EP[B, C]) extends Rep[B]
\end{lstlisting}
This is very easy, but looks inconsistent in usage. Another approach is to
not use a case class and provide a companion object with custom \cd{apply}
and \cd{unapply} methods. Listing~\ref{ligd-rtype} shows how this is
implemented.
\lstinputlisting[firstline=76,lastline=90,float=ht,caption=Implementation of RType,label=ligd-rtype]{scala/src/ligd.scala}

\begin{example}[Lists]
In order to represent a list, we encode the
list as a sum of products, or to be precise, we want to encode a \cd{List[A]}
as: \cd{Either[Unit, (A, List[A])]}. In order to do this, we need functions
that convert from \cd{List[A]} to \cd{Either[Unit, (A, List[A])]} and
vice versa. An implementation of those functions is given in figure~\ref{ligd-rlist}.
%-- We can then define an RType instance as:
%\lstinputlisting[firstline=104,lastline=108]{scala/src/ligd.scala}

\lstinputlisting[firstline=96,lastline=117,float=ht,caption=Implementation of rList,label=ligd-rlist]{scala/src/ligd.scala}

Note how the first parameter of \cd{RType} is recursive.

There is one issue with this representation though: It leads to deep recursion,
potentially causing a stack overflow in functions using it.
\end{example}

\subsection{Writing generic functions}
Writing generic functions using LIGD is simple. We just pass the object(s)
we are interested in to the function, along with their type representations,
and then pattern match on the type representations.

\begin{example}[Generic equality]
Listing~\ref{ligd-geq} shows how to implement a simple function that checks for
generic equality.

  \lstinputlisting[firstline=36,lastline=51,float=ht,caption=Generic Equality in LIGD,label=ligd-geq]{scala/src/ligd.scala}

This function does not deal with cycles nor does it support very long lists. In
the first case, it would need to keep track of which objects it visited already;
and in the second case, it would need to manage stack itself or be tail recursive.

Making the function tail-recursive seems straight forward. There is one issue,
though: Scala will not consider the calls to \cd{geq} eligible for tail call
optimization, as the type parameters differs. The only ways to deal with deeply
nested structures are thus manual stack management or trampolines\footnote{The
\cd{scala.util.control.TailCalls} object makes trampolines easy to write}. This
does not matter much in this chapter.
\end{example}


\subsection{Differences and Extensions to the original Haskell implementation}
One minor difference is that we implemented \cd{Rep} like we would implement
a type class in Haskell. The original paper manually passed around \cd{Rep}
instances for generic functions, we make use of Scala's implicit values.


\subsubsection{Ad-hoc cases using subtyping}
A key difference that extends the functionality of LIGD in Scala is that
Scala permits sub-typing. If we look at our definition of RType, we see
that it is not \cd{sealed}, allowing it to be extended by subclasses.

We can use this to create custom sub classes for types we are interested
in, and pattern match on them. Because those are sub classes of \cd{RType},
existing code will be unaffected by new types.

An example for this is the list type. Because we extended \cd{RType} as
\cd{RList} we can now write a generic function that finds all lists of a
specific element type in a data structure.

\subsection{Extensibility}
\label{lidg-extensible}
Imagine we have an existing generic function that we do not control (for example,
in a different module) and want to add a special case to that function. Because
the existing generic function has self-recursion this is usually impossible.

It turns out that if we change the way we write the generic function slightly,
we can actually extend it with new special cases. The changes are very simple:

\begin{enumerate}
  \item Convert the function into a class with an apply method containing the
        existing function and replace all recursive calls to the functions with
        calls to apply.
  \item Provide an object that extends that class, this can then be called like
        a function.
\end{enumerate}

We can now easily extend that `function', by extending the class and overriding
the apply method, matching on the cases we are interested in and deferring the
remaining cases to the implementation in the super class.

For example, if we transform \cd{geq} according to those rules, we can
easily extend it with a special case for lists, as shown in listing~\ref{src:ligd-extensible:geqlist}.\footnote{The complete example with the transformed \cd{geq} can be found in
the file~\texttt{src/ligd-extensible.scala} \cite{src}}
\lstinputlisting[firstline=76,lastline=85,float=ht,caption=Extended Generic Equality in LIGD,label=src:ligd-extensible:geqlist]{scala/src/ligd-extensible.scala}

A similar approach is used in `Lightweight modular staging'\cite{DBLP:conf/gpce/RompfO10}\cite{DBLP:journals/cacm/RompfO12},
although based on traits rather than classes, and not using an \cd{apply}
method -- thus making it easily possible to combine functions using inheritance
rather than composition of objects.

This approach using classes is not applicable to Haskell, but it is possible
to do something similar - by passing the function as an argument:
\begin{enumerate}
  \item Convert \cd{geq :: Rep[a] -> a -> a -> Bool} to \\
        \hphantom{Convert }\cd{geq' :: Rep[a] -> a -> a -> (forall a. Rep[a] -> a -> a -> Bool) -> Bool} \\
        and change self-recursive calls to call the new argument
  \item Add a  \cd{geq rep a b = geq' rep a b geq}
\end{enumerate}

Extending the function should now be straight forward. It is less useful to
have this in Haskell, though, because LIGD in Haskell does not support ad-hoc
cases, meaning we can only override the pre-defined type representations and not
match on custom types. If support for named constructors is added however, it
could still become useful enough.\footnote{See \url{http://code.haskell.org/generics/comparison/LIGD/LIGD.lhs} for an LIGD with named constructors}

\subsubsection{Equality comparisons of type representations}
The Haskell version of LIGD does not support comparing Rep instances
against each other. This is quite useful, though, as it allows to
implement simple variants of generic operations that look for objects
of a specific type in another object, such as folds.

For this to work, it is important to always use values for \cd{Rep} instances
for non-parameterized types, so that they compare correctly without a custom
\cd{equals()} method.
When representing parameterized types such as Lists, we can create a subclass
of RType with Rep parameters for the parameters of the type. We can then
define a custom \cd{equals()} method, as seen in listing~\ref{rmycustomtype}.

\begin{lstlisting}[float,caption=Pattern for representing parameterized types,label=rmycustomtype]
  class RMyCustomType[T1, ...](r1: Rep[T1], ...) {
    override def equals(b: Any) : Boolean = b match {
      case RMyCustomType(r1, ...) => r1 == this.r1 &&
                                     r2 == this.r2 && ...
      case _ => false
    }
  }
\end{lstlisting}

\section{EMGM -- Translation of Extensible and Modular Generics for the Masses}
\summary{More verbose (especially for ad-hoc cases) than LIGD}


In Generics for the Masses \cite{GM}, Ralf Hinze introduces an alternative
encoding of generic functions. While objects are still viewed as sums and
products of scalar values, generic functions are now defined as instances
of a \cd{Generic} type class, with one method for each type of object
supported, our slightly extended version can be seen in Figure~\ref{src:emgm:Generic}
\lstinputlisting[firstline=40,lastline=50,label=src:emgm:Generic,caption=The \cd{Generic} trait in GM]{scala/src/emgm.scala}
The parameter \cd{G} refers to a type constructor for a type that stores
a function of a specific type.

\begin{example}[Generic equality]
Imagine we want to check arbitrary objects of the same type for equality.
Our type \cd{G[\_]} must have a function that accepts two
value of the type given by the first parameter, and returns a boolean
-- see Listing~\ref{src:emgm:GEq}.
\lstinputlisting[firstline=172,lastline=192,float=ht,label=src:emgm:GEq,caption=A generic equality function in GM]{scala/src/emgm.scala}

We use closures here instead of normal methods in order to have the
compiler infer the types automatically, as explained in Section~\ref{type-inference}. \qed
\end{example}

Being able to represent generic functions is only half of the work, though
-- we still need a way to represent objects. Instead of converting objects
to a representation of sums and products, we simply call the members of
the \cd{Generic} instance we need.

\lstinputlisting[firstline=58,lastline=82,float=ht,label=src:emgm:Rep,caption=The \cd{Rep} trait in GM]{scala/src/emgm.scala}

We can now, similar to  LIGD, represent arbitrary types as sums of products.
Recall the list example from LIGD.
As we can see in Figure~\ref{src:emgm:rList}, the definition is similar to
the one in LIGD
-- the isomorphism is basically the same, the only difference is the slightly
more verbose definition of \cd{RList}, because we need to split the generic
function `creation' (the description of the type) out in \cd{rList}.

We will see that this only gets worse once we introduce ad-hoc cases.
\lstinputlisting[firstline=119,lastline=138,float=ht,label=src:emgm:rList,caption=Representing lists in GM]{scala/src/emgm.scala}

\subsection{Extensibility and Modularity}
Now consider that we want to treat lists specially, that is, create an
ad-hoc case of \cd{geq} that handles lists directly, instead of as a
sum of products representations. That is, we need a subclass of Generic
that can work on lists, something like:
\lstinputlisting[firstline=240,lastline=242,label=src:emgm:GenericList,caption=Generic with ad-hoc list case]{scala/src/emgm.scala}
Note that the special case for lists defers to the sum of product variant
if it is not handled by the generic function. We can thus use that extra
method (override it) or ignore it (and treat it using sums and products).
\lstinputlisting[firstline=253,lastline=270,float=ht,label=src:emgm:GEqList,caption=Generic equality specialised to lists]{scala/src/emgm.scala}

What we are still missing is a sort of dispatcher, like \cd{Rep} in Generics
for the masses. Without such a dispatcher, using generic functions like
\cd{GEqList} would not be easy.

\subsection{Comparison to LIGD}

\section{Uniplate}
\summary{Not feasible without macros or code generators and fails generic equality test}

Uniplate is a library that differentiates between two basic type of
objects: Uniplates and Biplates. Uniplates are simple recursive types
like terms or expressions, which form a graph where all inner nodes have
the same common supertype. Biplates are types that contain other Uniplates,
for example, a \cd{List a} has a \cd{Biplate List a}.

\subsection{Uniplates}
The \cd{Uniplate} type class, shown in listing~\ref{src:uniplate:Uniplate}
describes recursive types like terms and expressions. Its only primitive
operation is the \cd{uniplate()} method.

\lstinputlisting[firstline=38,lastline=56,label=src:uniplate:Uniplate,caption=The Uniplate type class (and various additional functions),style=breaklines]{scala/src/uniplate.scala}

The \cd{uniplate()} method takes an object of type \cd{T} and returns a tuple
of all direct children of the same type and a function that constructs a new \cd{T}
from a list of children of the same type.


There are also several more functions that can be defined on top of \cd{uniplate()}.

The \cd{children()} function returns a list of all direct children of the same
type. For example, in a term, this will be the operands.

The \cd{universe()} function returns a list of all transitive children of the
same type. For example, in a term, this will be the operands and their children,
and so on. In a tree, applied to the root, this will be all nodes in the tree.

The \cd{transform()} function defines a bottom up transformation. It applies
the passed function recursively to all children bottom-up.

\begin{example}[Transformation]
An example of a bottom-up transformation would be eliminating subtraction in expressions, that
is transform \cd{a - b} to \cd{a + (-b)}. Given an existing type \cd{Expr}
and \cd{Add}, \cd{Sub}, \cd{Neg} case classes of it, such a transformation
could be written as:

\begin{lstlisting}
    def removeSub(e: Expr) = {
        case Sub(a, b) => Add(a, Neg(b))
        case e => e
    }

    val expressionWithoutSub = transform(removeSub)(expression)
\end{lstlisting}
\end{example}

The \cd{descend()} function applies a top-down transformation.

The \cd{rewrite()} function applies a transformation until some sort of
normal form has been reached. The passed function returns a \cd{Some} if
it applied a transformation to an object and a \cd{None} if the input was
in normal form already.

In the original uniplate paper~\cite{DBLP:conf/haskell/MitchellR07}, Mitchell
and Runciman introduce additional functions, but those are not implemented
in the Scala version, because the focus was on biplates.

\subsection{Biplates}
Simply speaking, biplates are containers of uniplates. For example, a type
\cd{Expr} that represents arithmetic expressions over integers has a corresponding
\cd{Biplate[Expr, Int]}.

\lstinputlisting[firstline=99,lastline=103,label=src:uniplate:Biplate,caption=The Biplate type class (and various additional functions),style=breaklines]{scala/src/uniplate.scala}

The basic operation of a \cd{Biplate} is the \cd{biplate} method. This method
is similar to the \cd{uniplate} method of the \cd{Uniplate} type class, but now
works on children of a different type.

The \cd{universeBi} and \cd{transformBi} methods are like their non-bi counterparts
in \cd{Uniplate}. Outlook: \cd{transformBi} corresponds to the
\cd{everywhere} combinator of Scrap your Boilerplate~\cite{DBLP:conf/tldi/LammelJ03},
with one caveat: the passed function is not polymorphic.

\begin{example}[Minimum integer]
A function that calculates the minimum integer in a biplate can be defined
by folding \cd{scala.math.min} over the universe of the biplate.

\lstinputlisting[firstline=110,lastline=111,label=src:uniplate:minInt,caption=Uniplate minInt]{scala/src/uniplate.scala}
\end{example}

\subsection{Limitations and Advantages compared to other solutions}
Uniplate can theoretically be ported from Haskell to Scala, but it is not
feasible to do without a way to automatically generate instances of the
\cd{Biplate} type class, because the amount of biplates needed is quadratic
to the number of types.

Due to the way uniplate works, it is not possible to implement a generic
equality function for two biplate objects. Because biplates work by only
looking at a subset of the types elements (for example, a \cd{Biplate[Company,Salary]}
only cares about Salaries in a company -- not enough to compare two companies). Even
for Uniplates, defining a generic equality function won't work because we
only have access to the objects themselves and not to members of a different
type.

On the other hand, if Uniplate is applicable it will almost always lead to
the shortest code compared to other frameworks. See chapter~\ref{libraries}
for more examples.

\section{Shapeless -- Scala Implementation of Scrap your Boilerplate}

Shapeless is an adaption of Scrap your Boilerplate\cite{DBLP:conf/tldi/LammelJ03}
from Haskell to Scala. Compared to the Haskell implementation, the Scala
implementation simplifies some things a bit.

Comparing our ports of other libraries against shapeless helps us determine
whether the solutions are efficient.

I found Shapeless to be buggy. Defining a generic function and calling
it twice in the interpreter resulted in errors:
\begin{lstlisting}[language=Scala,gobble=2]
  scala> def inc(i: Int) = i + 1
  inc: (i: Int)Int

  scala> everywhere(inc _)(List(1,2,3))
  res7: List[Int] = List(2, 3, 4)

  scala> everywhere(inc _)(List(1,2,3))
  java.lang.AbstractMethodError: fresh$macro$6$2$.caseUniv()Lshapeless/PolyDefns$Case;
    ... 43 elided
\end{lstlisting}

Furthermore, testing it with a standard paradise benchmark, increasing all
salaries by 10\% failed as well sometimes, with shapeless only increasing
some of the values when run under scalatest.


    \chapter{Example Operations}

\section{Paradise benchmark}
GPBench, developed by Rodriguez et. al., is a suite of benchmarks for comparing
generic programming libraries in Haskell~\cite{DBLP:conf/haskell/RodriguezJJGKO08}.
One benchmark included in it is the so-called paradise benchmark.

The paradise benchmark looks at data structures representing a company
and tries to transform it, by increasing every salary by 10\%.{}

\lstinputlisting[firstline=28,lastline=39,float=ht,caption=Paradise benchmark data structures,label=src:data:company]{scala/src/data.scala}

The standard example data can be seen in Listing~\ref{src:data:company:example}.

\lstinputlisting[firstline=41,lastline=50,float=ht,caption=Paradise benchmark example data,label=src:data:company:example]{scala/src/data.scala}

Apart from increasing the salary, we will also look at querying salaries.


\subsection{LIGD}

Given the data structures above and our implementation of LIGD that supports
subtyping \cd{RType}, implementing both operations is fairly straight forward.

First of all, the type representations. For the data types \cd{Person},
\cd{Employee}, and \cd{Company} we can use the simplest form of definition,
by using the \cd{rType} function we added to our LIGD implementation.

\lstinputlisting[firstline=29,lastline=30,frame=t,caption={Person, Employee, and Company representations in LIGD},label=src:ligd-company:stuff,style=breaklines]{scala/src/ligd-company.scala}
\lstinputlisting[firstline=46,lastline=46,frame=b]{scala/src/ligd-company.scala}

Because we want to query and transform salaries, we need to treat salaries
specially, so we can pattern match on it (otherwise we would need to compare,
but pattern matching is easier to read):

\lstinputlisting[firstline=24,lastline=26,float=ht,caption=Salary representation in LIGD,label=src:ligd-company:salary]{scala/src/ligd-company.scala}

Finally, the \cd{DUnit} and \cd{Dept} types are mutually recursive. Simply
using \cd{rType} here as we have done for the other types would not work due
to two reasons: First, the Scala compiler would not be able to infer the types;
and second, because the implicit \cd{Rep} of \cd{rType} is not lazy, there would
be either endless recursion (if both were defined using \cd{def}) or \cd{rDept}
would be \cd{null} in the definition of \cd{rDUnit} because \cd{rDept} is defined
after \cd{rDUnit} and thus not yet initialised.

Fixing this is simple: For the first issue, we explicitly specify the types of
the representations. For the second issue, we can simply define \cd{rDUnit} in
terms of \cd{RType} instead of \cd{rType}. This makes the reference to \cd{rDept}
lazy and thus allows it to be used here.

\lstinputlisting[firstline=34,lastline=43,float=ht,caption=Mutually recursive representations,label=src:ligd-company:mutual]{scala/src/ligd-company.scala}

\begin{example}[Query: Summing all salaries]

Querying salaries is relatively straight forward. Pattern match on all standard
LIGD types, and recursive for the product, sum, RType types; but before matching
against \cd{RType} match against \cd{RSalary}:

\lstinputlisting[firstline=49,lastline=59,caption=Summing salaries in LIGD (manually),label=src:ligd-company:sumSalaryOld]{scala/src/ligd-company.scala}

We can, however, also use the generic fold we implemented earlier, the definition
then becomes a single line of code:

\lstinputlisting[firstline=70,lastline=70,caption=Summing salaries in LIGD (fold),label=src:ligd-company:sumSalary]{scala/src/ligd-company.scala}

Note: If we swapped around the first two parameters of \cd{gfoldl} we would not even
need to specify the \cd{Float} type, as long as we would pass \cd{0.0} instead
of 0.

\end{example}

\begin{example}[Transformation: Increase all salaries]

Increasing the salaries by 10\% is what is usually called the paradise
benchmark. Implementing this is straight forward as well in LIGD:

\lstinputlisting[firstline=61,lastline=68,caption=Increasing the salaries in LIGD,label=src:ligd-company:incSalary]{scala/src/ligd-company.scala}

We can abstract this by building an \cd{everywhere} function\footnote{a sample implementation of \cd{everywhere} is included in the source code} that applies
a transformation function, similar to EMGM and Uniplates \cd{transformBi} -- and
then calling it with a salary-transforming function on our company,
like this:
\begin{lstlisting}[caption=Increase the salaries using \cd{everywhere}]
  everywhere((i: Salary) ⇒ Salary(i.salary * 1.1F))(genCom)
\end{lstlisting}

\end{example}

In any case, both examples are easy to implement.


\subsection{EMGM}

Implementing the paradise benchmark in EMGM requires more boilerplate compared
to LIGD. This version is slightly more powerful than the LIGD variant (it
provides ad-hoc cases for all types in a company, compared to only \cd{Salary}),
but only because this makes the code more readable.

We first start by extending \cd{Generic} with new methods for all types of
company objects. This allows us to write functions with ad-hoc cases for all
of those types.

\lstinputlisting[firstline=38,lastline=58,caption=GenericCompany,label=src:emgm-company:GenericCompany,style=breaklines]{scala/src/emgm-company.scala}

Now we need some boilerplate to have implicit \cd{GRep} instances, as shown
in listing~\ref{src:emgm-company:GRep}.
\lstinputlisting[firstline=61,lastline=75,caption=Implicit instances of \cd{GRep},label=src:emgm-company:GRep]{scala/src/emgm-company.scala}

\begin{example}[Query: Summing all salaries]
Summing all salaries works as expected by using a very simple (although
verbose) definition:
\lstinputlisting[firstline=99,lastline=116,caption=Summing salaries in EMGM,label=src:emgm-company:sumSalary]{scala/src/emgm-company.scala}
\end{example}

\clearpage
\begin{example}[Transformation: Increase all salaries]
It's straight forward to implement this in EMGM, although the definition is
very verbose.
\lstinputlisting[firstline=78,lastline=96,caption=Using the transform pattern: Increase salary,label=src:emgm-company:incSalary]{scala/src/emgm-company.scala}
\end{example}

\subsection{Uniplate}
Uniplate is applicable here, and as promised in chapter 3, it has the shortest
solutions, due to its unique approach.\footnote{Not counting type representations}

\begin{example}[Query: Summing all salaries]
 Trivial to do in uniplate: Get all salaries and fold over them:
\begin{lstlisting}
  def sumSalary[C](self: C)(implicit bp: Biplate[C, Salary])
    = universeBi(self).foldLeft(0.0F)((f, s) => f + s.salary)
\end{lstlisting}
\end{example}

\begin{example}[Transformation: Increase all salaries]
 A transformation is done using the \cd{transformBi} function which takes a function of
 type \cd{A => A} and applies it to a value of a type \cd{B} using a \cd{Biplate[B,A]}.
\begin{lstlisting}
  transformBi((i: Salary) ⇒ Salary(i.salary * 1.1F))(genCom)
\end{lstlisting}
\end{example}


\subsection{Shapeless}

Using shapeless is much easier than using our self-created libraries because
there is no need to create type representations manually, shapeless takes care
of that for us.

\begin{example}[Query: Summing all salaries]
We can simply use the \cd{everything} combinator presented in chapter 3.

\lstinputlisting[firstline=75,lastline=83,caption=Summing salaries in Shapeless]{scala/src/shapeless.scala}

To recall, the \cd{everything} combinator consists of two parts: A traversal function,
here \cd{extractSalary}, that extracts the types we are interested in,
and a function to combine the results from the traversal, such as the
\cd{addFloat} here.
\end{example}

\begin{example}[Transformation: Increase all salaries]

Transformations like that are easy to write in Shapeless. All we need to do
is create a function that takes a \cd{Salary} and produces a new one and pass
that to Shapeless' \cd{everywhere} function.

\lstinputlisting[firstline=67,lastline=68,caption=Increasing salaries in Shapeless]{scala/tests/CompanyTest.scala}

\end{example}

\section{Locating the smallest integer in a datatype}
As another example, let's take a look at what it takes to find the smallest
integer in a data structure. In all cases, we return \cd{Int.MaxValue} if we cannot
find any (smaller) integers.

\paragraph{LIGD} The definition in LIGD is straight forward:
\lstinputlisting[firstline=288,lastline=295,caption=Minimum in LIGD,label=src:ligd:min]{scala/src/ligd.scala}

Alternatively, we can use one of the fold variants we implemented in our LIGD
port -- see the section~\ref{src:ligd} for more examples.

\paragraph{EMGM} In EMGM, the definition is relatively straight forward, similar
to LIGD, although with a bit more boilerplate:

\lstinputlisting[firstline=316,lastline=335,caption=Minimum in EMGM,label=src:emgm:gmin]{scala/src/emgm.scala}

\paragraph{Uniplate}

Finding the minimum in Uniplate could look something like this:

\lstinputlisting[firstline=110,lastline=111,caption=Minimum in Uniplate,label=src:emgm:gmin]{scala/src/uniplate.scala}

That is, you find all integers in the data structure recursively, and then
simply pick the minimum. As a special bonus, this only works on data structures
than can contain integers, because the \cd{Biplate} has \cd{Int} as a second
parameter.
\paragraph{Shapeless}
\todo{add shapeless min}

\section{Name-binding Operations}
\section{Derived Data Types}

    \chapter{Evaluation of the approaches}
\todo{Lots of stuff to finish}
\section{Library overview}
\begin{table}[ht]
  \begin{tabular}{c|cccc}
    Feature (and test)    & LIGD & EMGM & Uniplate & Shapeless \\
                 \hline
    Ad-hoc cases (company) & \checkmark & \checkmark & \checkmark & \checkmark \\
    Multi-parameter functions (geq) & \checkmark & \checkmark & ($\times$) & \checkmark
  \end{tabular}

  \caption{Feature overview}
\end{table}

\section{Performance of the libraries}

To test the performance of the libraries, we ran various tests with them. For
each test, we executed it 600 thousand times, and calculated the average
time of the last 300 thousand executions.

The exact testing framework and all performed tests can be seen in
section~\ref{src:test}.

\begin{table}[ht]
\begin{tabular}{c|r|r|r|r}
\bf{test}  & \bf{company}    & \bf{geq}        & \bf{min}        & \bf{sum}       \\
\hline
Direct     &          3,0 μs &          0,4 μs &          0,6 μs &          0,4 μs\\
Shapeless  &         39,3 μs &             N/A &          4,7 μs &          4,8 μs\\
LIGD       &         16,0 μs &         15,3 μs &         13,5 μs &         13,8 μs\\
EMGM       &         18,4 μs &         16,5 μs &         12,2 μs &         16,5 μs\\
\end{tabular}

\caption{Benchmark results}
\label{bench}
\end{table}

As we can see in in table~\ref{bench}, all tested libraries perform roughly
the same time-wise.

\section{Final remarks}

    \chapter{Conclusion}
Shapeless and LIGD are the best options for generic programming in Scala
right now. Shapeless provides a much larger feature set, and is also more
widely used, and should thus be preferred in production use.

The extended LIGD is now capable of supporting ad-hoc cases and extensibility,
making it more powerful then the original implementation in Haskell. Its use
of only basic language features and its easy to write representations and
functions make it an optimal choice for introducing generic programming to
students.

It would be interesting to look at automatic generation of LIGD representations
used macros; likewise, an implementation that represents constructors using
heterogeneous lists also seems a worthwhile idea, opening up new possibilities
like generic zippers. While an initial prototype of such a library exists in
the source code of this bachelor thesis\footnote{See \texttt{src/hligd.scala} and
\texttt{src/hlists.scala}\cite{src}}, it requires more work to make it work completely.

Another topic to look at with regards to LIGD is the handling of lists and
other container types: Representing them as deeply-nested pairs can cause
stack overflows when implementing algorithms in a naive way. Adding direct
support for sequences and/or implementing a set of combinators representing
commonly-needed operations using trampolines might be worthwhile.


    % Bibliography and lists of stuff
    \bibliographystyle{alpha}
    \bibliography{thesis}
    \lstlistoflistings

    The complete source code of the libraries and examples is publicly
    available at: \\
    \url{https://github.com/julian-klode/bsc-thesis-code/}
\end{document}
